\documentclass{article}
\usepackage{arxiv}
\usepackage[utf8]{inputenc}
\usepackage[T1]{fontenc}
\usepackage{hyperref}
\usepackage{url}
\usepackage{booktabs}
\usepackage{amsfonts}
\usepackage{amsmath}
\usepackage{amssymb}
\usepackage{nicefrac}
\usepackage{graphicx}
\usepackage{natbib}
\usepackage{doi}
\usepackage{tikz}
\usepackage{pgfplots}
\pgfplotsset{compat=1.18}
\usepgfplotslibrary{fillbetween}
\usetikzlibrary{shapes,arrows,positioning,calc,backgrounds,fit,patterns,decorations.pathreplacing}
\usepackage{amsthm}
\usepackage{algorithm}
\usepackage{algorithmic}
\usepackage{placeins}

\title{Prompt Degrees of Freedom in LLM-Based Social Science:\\ A Specification Curve Framework for Transparent Reporting}

\author{
  Jacob Crainic \\
  YCRG Management Sciences Lab \\
  \texttt{jacobcrainic2008@gmail.com}
}

\date{March 2026}

\begin{document}

\maketitle

\subsection*{Research Question}

When researchers use large language models as synthetic survey respondents, the prompt is an analytic choice.
Model selection, persona wording, question framing, system instructions, temperature, and few-shot exemplars each represent decision points that may shift substantive conclusions, yet current practice reports results from a single configuration and treats the rest as incidental.
This is the prompt-domain version of the ``garden of forking paths'' in psychology and economics \citep{gelman2014statistical, simonsohn2020specification}.
We introduce \textit{Prompt Specification Curve Analysis} (P-SCA): a method that maps how conclusions vary across the space of defensible prompt designs, decomposes outcome variance by prompt dimension, and produces a reporting standard for LLM-based social science.

\subsection*{Background}

\citet{argyle2023out} showed that demographically conditioned language models can approximate U.S. public opinion distributions, a paradigm they call ``silicon sampling.''
But fidelity has limits.
\citet{bisbee2024synthetic} found that it degrades for underrepresented populations.
\citet{sclar2024quantifying} documented that minor prompt perturbations (whitespace, label ordering) can shift NLP benchmark accuracy by up to 76 points.
\citet{karell2025integrating} identify prompt design as one of three core methodological challenges for generative AI in social science, and note that the field has no established approach to treating prompting as an experimental condition.

Prompt sensitivity is real, but the literature on it remains descriptive.
There is no method that lets researchers (1) enumerate defensible prompt designs for a given study, (2) measure how much each design dimension contributes to outcome variation, or (3) test whether a reported finding holds across the prompt multiverse or depends on a particular configuration.
Specification curve analysis \citep{simonsohn2020specification} solved this problem for traditional analytic choices; Steegen et al.\ \citep{steegen2016increasing} generalized it as ``multiverse analysis,'' in which every defensible analytic path is run and the distribution of outcomes reported.
Our \textit{prompt multiverse} borrows this architecture and applies it to the prompt design space.

\subsection*{Methodology}

We define a six-dimensional \textit{prompt multiverse} over the design choices in LLM-based opinion elicitation:
(1) model (GPT-5.4, GPT-5.4-nano, Claude Sonnet 4.6, Gemini 2.5 Flash),
(2) persona format (bare demographic tag, first-person narrative, third-person vignette, survey-style respondent block),
(3) question framing (direct elicitation, Likert scale, forced choice),
(4) system prompt register (neutral, academic researcher, survey administrator),
(5) sampling temperature (0.0, 0.3, 0.7, 1.0), and
(6) few-shot exemplar count (0, 1, 3).
Researchers make these choices routinely but rarely justify or report them.

We sample 600 specifications from this $4 \times 4 \times 3 \times 3 \times 4 \times 3 = 1{,}728$-cell space using Latin Hypercube Sampling \citep{mckay1979comparison}, covering the full multiverse without exhaustive enumeration.
Each specification is administered to 20 demographic profiles drawn from five battleground states (PA, MI, WI, AZ, GA), balanced by party, with 5 repeated draws per profile to estimate within-specification variance.

For each specification, we compute the partisan opinion gap (mean Democratic response minus mean Republican response) on three ANES items: government spending, immigration levels, and gun regulation.
If a finding holds in 95\% of specifications, it is likely a property of the underlying opinion distribution rather than an artifact of prompt construction.
If it holds in 55\%, the result depends on design choices that the researcher did not report.
We plot specification curves sorted by effect magnitude, following \citet{simonsohn2020specification}.

To identify which dimensions drive variation, we compute $\eta^2$ from one-way decompositions of response scores by each dimension, producing a rank-ordering of design choices by their contribution to outcome variation.
The full study extends this with Sobol sensitivity indices \citep{saltelli2008global} estimated via Saltelli's quasi-random scheme, which captures interaction effects between dimensions.
Joint inference follows the permutation-based procedure of \citet{simonsohn2020specification}, testing whether LLM partisan structure exceeds what a null model of prompt-driven noise would produce.

\subsection*{Data and Materials}

Benchmark items come from the American National Election Studies (ANES) 2020 and 2024 waves.
Demographic profiles are constructed from five-state battleground county data (Pennsylvania, Michigan, Wisconsin, Arizona, Georgia), with four profiles per state varying by party, age, education, race, and gender.
We scope the profile set to swing-state demographics, where LLM-based opinion research is most likely to be deployed for electoral and policy analysis.
All prompt templates, sampling code, API call logs, and analysis scripts will be released as an open-source pipeline.

\subsection*{Expected Contributions}

This study provides the first specification curve analysis of prompt design in a social science application.
Where prior work established that prompt sensitivity exists \citep{sclar2024quantifying}, we measure \textit{which} dimensions matter and \textit{how much} they matter.
The variance decomposition table, reporting sensitivity indices for each prompt dimension, gives researchers concrete guidance on where robustness checks are necessary and where they are not.

We also propose a \textit{Prompt Specification Reporting Standard}: authors document their prompt multiverse, include specification curve plots, and disclose the share of defensible specifications consistent with their claimed finding.
This responds to the transparency gap that \citet{karell2025integrating} identified, and gives reviewers a concrete criterion for evaluating LLM-based empirical work.

Because P-SCA treats model version as an explicit dimension, it also tracks how conclusions shift across model generations.
A study run on GPT-4o today and replicated on its successor next year can be compared within the same specification curve, making replication legible even as the underlying systems change.

The method is not specific to political opinion.
The same pipeline applies wherever LLMs serve as behavioral proxies: economic preference estimation, moral judgment elicitation, organizational survey simulation.
We release P-SCA as an open-source toolkit.

\bibliographystyle{plainnat}
\bibliography{references}

\end{document}
